\centerline{\bf \Huge Принцип Дирихле}


\begin{flushright}
        \scriptsize{— Ударишь меня с правой? Нет.\\
                    — Левой? Нет.\\
                    — Неужто ора-ора устроишь?!?\\
                    — ДА! ДА! ДА! ДА!}
\end{flushright}


\problem{0.1}
Можно ли рассадить двух Баиров и Влада за одну парту?

\problem{0.2}
Можно ли разложить 44 шарика на 9 кучек так, чтобы количество шариков во всех кучках было различным (и ненулевым)?

\problem{1}
В турнире по Brawl Stars учавствовали 20 игроков. Самому старшему из них 35 лет, а самому младшему а) 16 лет б) 17 лет. Верно ли, что среди игроков есть одногодки?
% а) да, б) нет

\problem{2}
В ящике лежат шары: 5 красных, 7 синих и 1 зеленый. Сколько шаров надо вынуть, не глядя, чтобы наверняка достать 2 шара одного цвета?
% 4


\problem{3}
Олимпиаду в ЭлМШ писали 70 школьников. Составитель набрал максимальные 33 балла, остальные меньше. Докажите, что по крайней мере три школьника набрали одинаковое количество баллов.
% 33*3 = 69
\newline

\centerline{\bf \Huge Разнобой}

\problem{4}
Из 100 ребят, отправляющихся в детский оздоровительный лагерь, кататься на сноуборде умеют 30 ребят, на скейтборде — 28, на роликах — 42. На скейтборде и на сноуборде умеют кататься 8 ребят, на скейтборде и на роликах — 10, на сноуборде и на роликах — 5, а на всех трех — 3. Сколько ребят не умеют кататься ни на сноуборде, ни на скейтборде, ни на роликах?
% 20

\problem{5}
В комнате 10 островитян. Один из них сказал: "Здесь нет ни одного рыцаря"; второй: "Здесь не более одного рыцаря"; третий: "Здесь не более двух рыцарей", и т.д., десятый: "Здесь не более девяти рыцарей". Сколько в комнате рыцарей?
% 1

\problem{6}
Можно ли таблицу 5×5 заполнить числами -1, 0, 1 так, чтобы суммы во всех строках, во всех столбцах и на главных диагоналях были различны?
% нельзя 

\problem{7}
Среди 9 монет 2 фальшивые. Определите фальшивые монеты за 4 взвешивания на чашечных весах без гирь, если известно, что обе фальшивые монеты весят одинаково, причем тяжелее настоящих.